\pdfvariable suppressoptionalinfo 611
\pdfvariable trailerid {[<C2492026083000000000000000000000><C2492026083000000000000000000000>]}
\providecommand{\CRevisionRound}{2}
\documentclass[10pt]{article}
\usepackage[a4paper,margin=17mm]{geometry}
\usepackage{amsmath,amssymb,amsthm,booktabs,microtype}
\usepackage[T1]{fontenc}
\usepackage{lmodern}
\usepackage[hidelinks,hypertexnames=false]{hyperref}
\emergencystretch=2em
\newtheorem{theorem}{Theorem}
\newtheorem{proposition}[theorem]{Proposition}
\newcommand{\dd}{\,\mathrm d}
\newcommand{\R}{\mathbb R}
\newcommand{\VdP}{Van der Pol}
\title{A Sign-Complete Lienard Certificate for the Van der Pol Flow}
\author{HCS Research Program}
\date{30 August 2026\quad (revision \CRevisionRound)}
\begin{document}
\maketitle
\begin{abstract}
We freeze the smooth polynomial \VdP\ oscillator and take one theorem-scale
step beyond a local phase portrait.  For $\mu>0$, its Lienard primitive has a single
positive zero, so the classical theorem gives exactly one hyperbolic attracting
limit cycle for positive damping; time reversal gives one repelling cycle for
negative damping, while the zero face is a harmonic center with a continuum
of ovals.  Energy balance and planar divergence then identify the Poincare
fixed section and the transverse Floquet multiplier.  Five independently
replayed DOP853 rows illustrate the theorem (the multiplier decreases from
\(0.5331\) at \(\mu=0.1\) to \(1.28\times10^{-25}\) at \(\mu=4\)).
The rows are finite regression probes, not an all-state census.  This
source-local result makes no arithmetic, target-determinant, or
Hilbert--Polya claim.
\end{abstract}

\section{Frozen model and question}
We study the polynomial vector field on \(\R^2\)
\begin{equation}
 \dot x=y,\qquad \dot y=\mu(1-x^2)y-x,\qquad \mu\in\R .
 \label{eq:vdp}
\end{equation}
Equivalently, \(x\) satisfies
\begin{equation}
 x''+\mu(x^2-1)x'+x=0 .
 \label{eq:lienard}
\end{equation}
The clock is physical time \(t\), and the oriented Poincare section is
\(\Sigma=\{(x,y):x=0,\ y>0\}\).  The question is whether this smooth
nonlinear family can support a complete, replayable sign and boundary
statement without importing an arithmetic carrier.  The answer below is
analytic; the numerical receipt only supplies a transparent finite check.

\ifnum\CRevisionRound=0
The baseline freezes (\ref{eq:vdp})--(\ref{eq:lienard}), the section
orientation, and the three parameter faces \(\mu<0\), \(\mu=0\), and
\(\mu>0\).  No period is assigned by sampling alone.
\fi

\section{The Lienard theorem and sign boundary}
Write \(f(x)=\mu(x^2-1)\) and
\begin{equation}
 F(x)=\int_0^x f(s)\dd s=\mu\left(\frac{x^3}{3}-x\right).
 \label{eq:F}
\end{equation}
For \(\mu>0\), \(F<0\) on \((0,\sqrt3)\), has the unique positive zero
\(\sqrt3\), and is positive and increasing thereafter.  The standard
Liénard existence and uniqueness hypotheses therefore apply.

\begin{theorem}[sign-complete cycle statement]
For every \(\mu>0\), the flow (\ref{eq:vdp}) has exactly one hyperbolic
attracting periodic orbit surrounding the origin.  Its Poincare return map on
\(\Sigma\) has one fixed point, and the derivative at that point is the
transverse Floquet multiplier.  For \(\mu<0\), the map
\((t,x,y)\mapsto(-t,x,-y)\) transfers this orbit to exactly one hyperbolic
repelling cycle.  At \(\mu=0\),
\(E=(x^2+y^2)/2\) is conserved and every level \(E=c>0\) is a harmonic oval
of period \(2\pi\).
\end{theorem}

\paragraph{Proof sketch.}
Equation (\ref{eq:lienard}) is in the classical Liénard form.  The sign and
monotonicity of \(F\) in (\ref{eq:F}), together with the odd symmetry, give
existence and uniqueness of the enclosing cycle; hyperbolicity and attraction
are part of the same return-map conclusion.  Reversing time and velocity
changes the sign of \(\mu\), hence reverses stability.  Setting \(\mu=0\)
reduces (\ref{eq:vdp}) to the unit harmonic oscillator, proving the center
face directly.  This invokes no numerical fit and excludes no state by a
finite cutoff.

\ifnum\CRevisionRound>0
\section{Energy, divergence, and Floquet data}
The elementary identities are
\begin{equation}
 E=\frac{x^2+y^2}{2},\qquad \dot E=\mu(1-x^2)y^2,\qquad
 \operatorname{div}X=\mu(1-x^2).
 \label{eq:identities}
\end{equation}
For a periodic orbit with \(\mu\ne0\), integrating the first identity gives
the exact balance \(\int_0^T(1-x^2)y^2\dd t=0\); the divided balance is not
asserted on the \(\mu=0\) center face.  In two dimensions one Floquet multiplier is
one (the tangent direction), and Liouville's formula gives
\begin{equation}
 \lambda_\perp=\exp\!\left(\int_0^T\operatorname{div}X\dd t\right)
 =\exp\!\left(\mu\int_0^T(1-x^2)\dd t\right).
 \label{eq:floquet}
\end{equation}
For a restoring frequency \(\omega>0\), the change \(\tau=\omega t\)
reduces \(x''+\mu_0(x^2-1)x'+\omega^2x=0\) to the same normal form with
effective parameter \(\mu_0/\omega\).
\fi

\ifnum\CRevisionRound>0
\section{Finite Poincare receipt}
The producer brackets the fixed point by \(y\in[0.05,8]\) and integrates
with DOP853 (relative tolerance \(3\times10^{-12}\), absolute tolerance
\(3\times10^{-14}\), maximum step \(0.03\)).  The independent checker repeats
the integration and verifies the return residual and (\ref{eq:identities})--
(\ref{eq:floquet}).  Table~\ref{tab:receipt} reports rounded values; all
digits in the JSON receipt are retained at 15 significant figures.
\begin{table}[t]
\centering\small
\setlength{\tabcolsep}{4pt}
\caption{Positive-parameter return probes.  These rows are finite regression
receipts, not a numerical census or a proof of uniqueness.}
\label{tab:receipt}
\begin{tabular}{c|c|c|c|c|c}
\toprule
\(\mu\) & \(y_\Sigma\) & \(T\) & \(\int\!\operatorname{div}X\) & \(\lambda_\perp\) & |balance|\\
\midrule
0.1 & 2.00177 & 6.28711 & $-0.62910$ & $5.33\times10^{-1}$ & $6.2\times10^{-15}$\\
0.5 & 2.04406 & 6.38068 & $-3.23967$ & $3.92\times10^{-2}$ & $2.1\times10^{-13}$\\
1   & 2.17271 & 6.66329 & $-7.05893$ & $8.60\times10^{-4}$ & $1.8\times10^{-13}$\\
2   & 2.61497 & 7.62987 & $-18.17864$ & $1.27\times10^{-8}$ & $8.3\times10^{-13}$\\
4   & 3.76234 & 10.20352 & $-57.32131$ & $1.28\times10^{-25}$ & $5.4\times10^{-12}$\\
\bottomrule
\end{tabular}
\end{table}
\fi

\ifnum\CRevisionRound=1
\paragraph{Independent checks.}
The producer-independent checker passes 264 structural and numerical
assertions.  A separate SymPy reconstruction passes 81 identities, including
the origin characteristic polynomial, time reversal, and frequency scaling.
Clean-process byte replay succeeds, and 40 repaired-hash hostile mutations
are all rejected.  These gates validate the receipt's provenance and algebra;
they do not enlarge the theorem's scope.
\fi

\ifnum\CRevisionRound>1
\section{Collision audit and route boundary}
The model is not a relabeling of nearby packages.  C227 studies the
three-dimensional Lorenz flow and local equilibrium/Hopf conditions; C232
studies a conservative Duffing separatrix; C178 uses a harmonic strobe; C237
is a stochastic Kramers system; and C245 is hybrid integrate-and-fire.  C249
is the smooth polynomial Liénard case whose single-cycle theorem and Floquet
receipt are absent from those certificates.  This comparison is workspace
bookkeeping, not a literature-priority claim.

The locked scope is \texttt{NO\_BAD\_EULER\_OR\_ROOT\_NUMBER}.  No prime or
zero table, arithmetic local datum, Euler factor, root number, automorphy
claim, target divisor or functional equation, target determinant, or
Hilbert--Pólya operator is present.  Accordingly the Route-A tuple is
\texttt{(A0\_FAIL,A1\_PASS\_ANALYTIC,A2\_FAIL,A3\_FAIL,A4\_FORMAL\_HINT)},
with overall verdict \texttt{ROUTE\_A\_REJECTED}; Route B is disabled.
\fi

\section{Conclusion and limitations}
The main result is a single sign-complete theorem package: Liénard uniqueness
for positive damping, time reversal for negative damping, an exact center
boundary, and a divergence-to-Floquet identity tied to one Poincare section.
The receipt is deliberately modest.  It does not give an elementary formula
for the period at every parameter, enumerate every continuum state, or infer
asymptotic coefficients from five samples.  A future analytic direction is a
uniform asymptotic study of the relaxation period; that direction would still
be a smooth-flow problem and would not authorize arithmetic claims.

\begin{thebibliography}{9}
\bibitem{Lienard1928} A.~Liénard, ``Étude des oscillations entretenues,''
\emph{Revue Générale de l'Électricité} 23 (1928), 901--912,
\href{https://gallica.bnf.fr/ark:/12148/bpt6k5671115f}{Gallica scan}.
\bibitem{LevinsonSmith1942} N.~Levinson and O.~K.~Smith, ``A general equation
for relaxation oscillations,'' \emph{Duke Mathematical Journal} 9 (1942),
382--403, DOI
\href{https://doi.org/10.1215/S0012-7094-42-00928-1}{10.1215/S0012-7094-42-00928-1}.
\bibitem{Strogatz2015} S.~H.~Strogatz, \emph{Nonlinear Dynamics and Chaos},
2nd ed., Westview Press, 2015.
\end{thebibliography}
\end{document}
